\chapter{Skrip Visualisasi dan Pembangkitan Grafik Hasil Simulasi}

Pada bagian ini, dilampirkan kumpulan skrip pemrograman \textit{Python} yang secara khusus dikembangkan untuk mengeksekusi operasi pembuatan visualisasi data kuantitatif dari hasil simulasi \textit{Variational Quantum Eigensolver} (VQE) dan ekuilibrium teori permainan. Skrip-skrip ini berfungsi sebagai instrumen konversi dari representasi data komputasional mentah menjadi bentuk diagram grafis yang komprehensif, mencakup pembentukan kurva konvergensi energi historis, histogram sebaran probabilitas status kuantum, grafik perbandingan imbal hasil portofolio, hingga komposisi metrik rasio absolut Sharpe. Rangkaian rutinitas komputasional visual ini merupakan tulang punggung utama dari analisis observabel data hasil evaluasi performa model yang memfasilitasi pelacakan konvergensi lanskap energi portofolio lintas dimensi aset.

\begin{lstlisting}[language=Python, caption={Skrip utama generator grafik evaluasi portofolio (\texttt{plot\_generator.py})}, label={lst:plot-generator}]
import matplotlib
matplotlib.use('Agg')  # Menggunakan backend non-interaktif
import matplotlib.pyplot as plt
import numpy as np
import pennylane as qml
import pandas as pd
import os

# Import Qiskit untuk perenderan rangkaian yang lebih baik
from qiskit import QuantumCircuit, QuantumRegister, ClassicalRegister

def plot_all(results, data_clean, value_benchmark_idx, config, filenames):
    """
    Menghasilkan semua grafik yang diperlukan untuk analisis.
    """
    tickers = config['tickers']
    N = len(tickers)
    benchmark_ticker = config['benchmark_ticker']

    # 1. Folder Khusus Analisis Jendela
    output_dir = f"Analisis_Window_N{N}"
    if not os.path.exists(output_dir):
        os.makedirs(output_dir)
    
    print(f"\n--- Menghasilkan Grafik Analisis Per Jendela di '{output_dir}' ---")
    for i, win_data in enumerate(results['window_analysis_history']):
        win_filename = os.path.join(output_dir, f"{win_data['date']}_window.png")
        plot_window_analysis(win_data, n_qubits=N, filename=win_filename)
        
        # Render sirkuit per jendela
        circ_filename = os.path.join(output_dir, f"{win_data['date']}_circuit.png")
        plot_quantum_circuit(N, win_data['best_depth'], circ_filename, params=win_data['best_params'])

    # 2. Update nama file rangkaian untuk depth terakhir (untuk folder Hasil_NX)
    final_depth = results['depths_history'][-1] if len(results['depths_history']) > 0 else 1
    suffix = f"_N{N}"
    filenames['circuit_png'] = f'rangkaian_kuantum_depth{final_depth}{suffix}.png'
    
    # Rangkaian Kuantum Window Terakhir dengan Parameter
    final_win = results['window_analysis_history'][-1]
    plot_quantum_circuit(N, final_depth, filenames['circuit_png'], params=final_win['best_params'])

    # 4. Hasil Backtest VQE vs Benchmarks
    # Cari data classic jika ada
    classic_file = f"classic_compare/Hasil_Classic_Compare/equity_history_classic_N{N}.csv"
    value_classic = None
    if os.path.exists(classic_file):
        try:
            df_classic = pd.read_csv(classic_file)
            value_classic = df_classic['Equity'].values
        except:
            pass

    plot_equity_growth(data_clean, results['value_vqe'], results['value_bench'],
                       value_benchmark_idx, benchmark_ticker, 'Quantum VQE',
                       'blue', filenames['backtest_vqe_png'], value_classic=value_classic)

    # 5. Hasil Backtest Nash vs Benchmarks
    plot_equity_growth(data_clean, results['value_nash'], results['value_bench'],
                       value_benchmark_idx, benchmark_ticker, 'Nash Equilibrium (Klasik)',
                       'orange', filenames['backtest_nash_png'], value_classic=value_classic)

    # 6. Depth per Window
    plot_depth_per_window(results['rebalance_dates'], results['depths_history'], filenames['depth_png'])

def plot_window_analysis(win_data, n_qubits, filename):
    """
    Merender 6 panel analisis untuk satu jendela waktu tertentu.
    """
    date = win_data['date']
    best_history = win_data['best_history']
    best_ent_hist = win_data['best_ent_hist']
    depth_energies = win_data['depth_energies']
    lr_data = win_data['lr_data']
    probs = win_data['winning_probs']
    K = win_data['K']
    use_warm_start = win_data.get('use_warm_start', True)

    # Hitung Brute Force khusus untuk jendela ini (dengan suffix N)
    bf_energy = calculate_bf_for_window(date, n_qubits)

    fig, axes = plt.subplots(3, 2, figsize=(16, 15))
    
    # Judul Dinamis berdasarkan Warm-Start
    ws_text = "Warm-Start" if use_warm_start else "tanpa Warm-Start"
    fig.suptitle(f"Analisis VQE {ws_text} - Jendela {date} (N={n_qubits}, K={K})", fontsize=20, fontweight='bold')
    axes = axes.flatten()

    # 1. Konvergensi Energi
    iters = np.arange(1, len(best_history) + 1) * 10 # batch_size=10
    final_e_spsa = best_history[-1]
    axes[0].plot(iters, best_history, marker='o', markersize=3, color='blue', label=f'SPSA Progress ({final_e_spsa:.4f})')
    if bf_energy is not None:
        axes[0].axhline(y=bf_energy, color='red', linestyle='--', label=f'Brute Force ({bf_energy:.4f})')
        all_e = list(best_history) + [bf_energy]
        axes[0].set_ylim(min(all_e)-0.01, max(all_e)+0.01)
    axes[0].set_title('1. Konvergensi Energi (Iterasi SPSA)')
    axes[0].set_ylabel('Energi')
    axes[0].legend()
    axes[0].grid(True, alpha=0.3)

    # 2. Energi vs Depth
    d_list = [d for d, e, it in depth_energies]
    e_list = [e for d, e, it in depth_energies]
    vqe_min_depth = min(e_list)
    axes[1].plot(d_list, e_list, marker='s', markersize=8, color='blue', label=f'VQE Min ({vqe_min_depth:.4f})')
    if bf_energy is not None:
        axes[1].axhline(y=bf_energy, color='red', linestyle='--', label=f'Brute Force ({bf_energy:.4f})')
        all_e = e_list + [bf_energy]
        axes[1].set_ylim(min(all_e)-0.01, max(all_e)+0.01)
    axes[1].set_title('2. Pencarian Energi vs Depth')
    axes[1].set_xticks(d_list)
    axes[1].legend()
    axes[1].grid(True, alpha=0.3)

    # 3. Entanglement Entropy
    if len(best_ent_hist) > 0:
        axes[2].plot(iters, best_ent_hist, color='purple', marker='^', markersize=3)
        axes[2].set_title('3. Entanglement Entropy vs Iterasi')
        axes[2].set_ylabel('Von Neumann Entropy')
        axes[2].grid(True, alpha=0.3)

    # 4. Learning Rate Finder
    test_a, test_e, best_a = lr_data
    axes[3].plot(test_a, test_e, marker='o', color='green')
    axes[3].axvline(best_a, color='red', linestyle='--', label=f'Best a: {best_a:.4f}')
    axes[3].set_xscale('log')
    axes[3].set_title('4. LR Finder (Calibration)')
    axes[3].legend()
    axes[3].grid(True, alpha=0.3)

    # 5. Distribusi Probabilitas (Bar Chart)
    bitstrings = [format(i, f'0{n_qubits}b') for i in range(2**n_qubits)]
    colors = ['limegreen' if b.count('1') == K else 'salmon' for b in bitstrings]
    axes[4].bar(bitstrings, probs, color=colors, edgecolor='black')
    axes[4].set_title(f'5. Probabilitas State (K={K})')
    axes[4].set_ylim(0, 1.1)
    # Tambahkan label teks di atas batang
    for idx, p in enumerate(probs):
        axes[4].text(idx, p + 0.02, f"{p:.2f}", ha='center', fontsize=8)

    # 6. Variansi Gradien (Diagnostic untuk Barren Plateaus)
    best_gvar_hist = win_data.get('best_gvar_hist', [])
    if len(best_gvar_hist) > 0:
        axes[5].plot(iters, best_gvar_hist, color='orange', marker='x', markersize=3)
        axes[5].set_yscale('log') # Gunakan skala log karena variansi bisa sangat kecil
        axes[5].set_title('6. Variansi Gradien vs Iterasi (Log)')
        axes[5].set_ylabel('Var(Grad)')
        axes[5].grid(True, alpha=0.3)
    else:
        axes[5].text(0.5, 0.5, 'Data Gradien Tidak Tersedia', ha='center')

    plt.tight_layout(rect=[0, 0.03, 1, 0.95])
    plt.savefig(filename)
    plt.close()

def plot_quantum_circuit(N, depth, filename, params=None, title=None):
    """
    Merender rangkaian kuantum menggunakan Qiskit untuk hasil yang lebih berwarna,
    informatif (label parameter), dan otomatis dibagi baris (folding).
    """
    qr = QuantumRegister(N, 'q')
    cr = ClassicalRegister(N, 'c')
    qc = QuantumCircuit(qr, cr)

    if params is None:
        params = np.zeros(N * 2 * (depth + 1))
    
    # Reshape parameter sesuai struktur Ansatz (layer, qubit, gate_type)
    w = params.reshape((depth + 1, N, 2))
    
    for layer in range(depth + 1):
        for q_idx in range(N):
            # Ambil nilai theta dan bulatkan agar tidak terlalu panjang di gambar
            theta_ry = np.round(float(w[layer, q_idx, 0]), 2)
            theta_rz = np.round(float(w[layer, q_idx, 1]), 2)
            
            # Tambahkan gate RY dan RZ dengan label parameter numerik
            qc.ry(theta_ry, qr[q_idx])
            qc.rz(theta_rz, qr[q_idx])
        
        if layer < depth:
            # Tambahkan Barrier antar layer untuk kejelasan visual
            qc.barrier()
            # Chain CNOT (Entanglement)
            for q_idx in range(N - 1):
                qc.cx(qr[q_idx], qr[q_idx + 1])
            qc.cx(qr[N - 1], qr[0])
            qc.barrier()

    # Tambahkan Pengukuran di akhir
    qc.measure(qr, cr)

    # Pengaturan Visual Qiskit
    # style='iqp' memberikan skema warna modern/IBM standard
    # fold=20 disetel agar muat sekitar 5 depth per baris (mencegah terlalu panjang ke samping)
    fig = qc.draw(output='mpl', style='iqp', plot_barriers=True, justify='left', fold=20)
    
    if title:
        fig.suptitle(title, fontsize=16, fontweight='bold', y=1.05)
    else:
        fig.suptitle(f"VQE Optimized Circuit (N={N}, Depth={depth})", fontsize=16, fontweight='bold', y=1.05)
        
    fig.savefig(filename, bbox_inches='tight', dpi=150)
    plt.close(fig)
    print(f"Gambar rangkaian kuantum (Qiskit) disimpan sebagai '{filename}'.")

def calculate_bf_for_window(date, n_qubits):
    """Fungsi pembantu untuk hitung BF on-the-fly dengan Penalty Annealing Target."""
    suffix = f"_N{n_qubits}"
    try:
        # 1. Ambil Bias (h)
        h_df = pd.read_csv(f'bias_h_total{suffix}.csv')
        h_win = h_df[h_df['Date'] == str(date)]
        h_obj = h_win['Bias_h_Obj'].values
        h_pen = h_win['Bias_h_Pen'].values
        n = len(h_obj)
        
        # 2. Ambil Interaksi (J)
        J_df = pd.read_csv(f'interaksi_J_total{suffix}.csv')
        J_win = J_df[J_df['Date'] == str(date)]
        J_mat_obj = np.zeros((n, n))
        J_mat_pen = np.zeros((n, n))
        tickers = h_win['Ticker'].values
        t2i = {t: i for i, t in enumerate(tickers)}
        for _, r in J_win.iterrows():
            i, j = t2i[r['Ticker_i']], t2i[r['Ticker_j']]
            J_mat_obj[i,j] = J_mat_obj[j,i] = r['Interaction_J_Obj']
            J_mat_pen[i,j] = J_mat_pen[j,i] = r['Interaction_J_Pen']
            
        # 3. Ambil Konstanta C dan Target Penalty
        c_df = pd.read_csv(f'parameter_pendamping{suffix}.csv')
        c_win = c_df[c_df['Date'] == str(date)].iloc[0]
        c_obj = c_win['C_Obj']
        c_pen = c_win['C_Pen']
        penalty_A = c_win['Penalty_A_Target']
        
        # 4. Kombinasikan menjadi Hamiltonian Total Akhir
        h_total = h_obj + penalty_A * h_pen
        J_total = J_mat_obj + penalty_A * J_mat_pen
        c_total = c_obj + penalty_A * c_pen
        
        from itertools import product
        min_e = float('inf')
        for b in product([0, 1], repeat=n):
            Z = 1 - 2 * np.array(b)
            # E = sum(h_i * Z_i) + sum(J_ij * Z_i * Z_j) + C
            e = np.sum(h_total * Z) + sum(J_total[i,j]*Z[i]*Z[j] for i in range(n) for j in range(i+1, n)) + c_total
            if e < min_e: min_e = e
        return min_e
    except:
        return None

    plt.tight_layout()
    plt.savefig(filename)
    print(f"Grafik konvergensi detail disimpan sebagai '{filename}'.")

def plot_equity_growth(data_clean, value_strategy, value_bench, value_benchmark_idx, benchmark_ticker, strategy_label, color, filename, value_classic=None):
    plt.figure(figsize=(12, 6))
    plt.plot(data_clean.index[:len(value_strategy)], value_strategy, label=strategy_label, linewidth=2.5, color=color)
    plt.plot(data_clean.index[:len(value_bench)], value_bench, label='Benchmark (Equal Weight)', linestyle=':', color='black', alpha=0.6)
    plt.plot(value_benchmark_idx.index, value_benchmark_idx.values, label=f'Indeks ({benchmark_ticker})', linestyle='-.', color='magenta', linewidth=2)
    
    if value_classic is not None:
        plt.plot(data_clean.index[:len(value_classic)], value_classic, label='Classic Markowitz (Mean-Var)', linestyle='--', color='green', linewidth=2)

    plt.title(f'Perbandingan Pertumbuhan Ekuitas: {strategy_label} vs Benchmarks')
    plt.ylabel('Total Ekuitas (Rupiah)')
    plt.xlabel('Tanggal')
    plt.legend(loc='upper left')
    plt.grid(True, linestyle=':', alpha=0.7)
    plt.tight_layout()
    plt.savefig(filename)
    print(f"Grafik {strategy_label} disimpan sebagai '{filename}'.")

def plot_depth_per_window(rebalance_dates, depths_history, filename):
    plt.figure(figsize=(10, 5))
    avg_depth_floor = int(np.floor(np.mean(depths_history)))
    plt.plot(rebalance_dates, depths_history, marker='o', linestyle='-', color='purple', label='Depth Terpilih')
    plt.axhline(y=avg_depth_floor, color='red', linestyle='--', label=f'Rata-rata Depth (Floor): {avg_depth_floor}')
    plt.title('Depth Terpilih vs Rebalance Window')
    plt.ylabel('Depth')
    plt.xlabel('Tanggal Rebalance')
    plt.xticks(rotation=45)
    plt.legend()
    plt.grid(True, linestyle=':', alpha=0.7)
    plt.tight_layout()
    plt.savefig(filename)
    print(f"Grafik depth per window disimpan sebagai '{filename}'. Rata-rata (floor): {avg_depth_floor}")
\end{lstlisting}

\begin{lstlisting}[language=Python, caption={Skrip visualisasi distribusi probabilitas status VQE (\texttt{render\_vqe\_probs.py})}, label={lst:render-vqe-probs}]
import matplotlib
matplotlib.use('Agg')
import matplotlib.pyplot as plt
import numpy as np
import pennylane as qml
import pandas as pd
import os

def render_vqe_probability_bar(n_qubits, K):
    # 1. Ambil Parameter Terbaik (Theta) dari hasil run sebelumnya
    # Kita asumsikan file 'theta_final_all_depths.csv' ada
    if not os.path.exists('theta_final_all_depths.csv'):
        print("Error: File theta_final_all_depths.csv tidak ditemukan. Jalankan simulasi terlebih dahulu.")
        return

    df_theta = pd.read_csv('theta_final_all_depths.csv')
    
    # Ambil depth terakhir (max depth) yang tercatat
    max_depth_in_file = df_theta['Depth'].max()
    params = df_theta[df_theta['Depth'] == max_depth_in_file]['Theta_Value'].values
    
    print(f"Merender Probabilitas untuk Depth: {max_depth_in_file}")

    # 2. Definisikan Sirkuit Kuantum untuk mendapatkan Probabilitas
    dev = qml.device("default.qubit", wires=n_qubits)

    @qml.qnode(dev)
    def prob_circuit(p):
        depth = max_depth_in_file
        w = p.reshape((depth + 1, n_qubits, 2))
        for layer in range(depth + 1):
            for q in range(n_qubits):
                qml.RY(w[layer, q, 0], wires=q)
                qml.RZ(w[layer, q, 1], wires=q)
            if layer < depth:
                for q in range(n_qubits - 1):
                    qml.CNOT(wires=[q, q + 1])
                qml.CNOT(wires=[n_qubits - 1, 0])
        return qml.probs(wires=range(n_qubits))

    # Hitung Probabilitas
    probs = prob_circuit(params)
    bitstrings = [format(i, f'0{n_qubits}b') for i in range(2**n_qubits)]

    # 3. Visualisasi Grafik Batang
    plt.figure(figsize=(10, 6))
    
    # Tentukan warna: Hijau untuk yang valid (sum=K), Merah untuk yang melanggar
    colors = []
    for bs in bitstrings:
        if bs.count('1') == K:
            colors.append('limegreen')
        else:
            colors.append('salmon')

    bars = plt.bar(bitstrings, probs, color=colors, edgecolor='black', alpha=0.8)
    
    # Tambahkan Label Probabilitas di atas batang
    for bar in bars:
        yval = bar.get_height()
        plt.text(bar.get_x() + bar.get_width()/2, yval + 0.01, f'{yval:.3f}', ha='center', va='bottom', fontsize=10)

    # Tambahkan keterangan/legend manual
    from matplotlib.lines import Line2D
    legend_elements = [
        Line2D([0], [0], color='limegreen', lw=4, label=f'Valid (K={K})'),
        Line2D([0], [0], color='salmon', lw=4, label='Invalid (Penalty Applied)')
    ]
    plt.legend(handles=legend_elements, loc='upper right')

    plt.title(f"Distribusi Probabilitas State VQE (N={n_qubits}, K={K})", fontsize=14)
    plt.ylabel("Probabilitas", fontsize=12)
    plt.xlabel("Bitstring (Spin Interpretation)", fontsize=12)
    plt.ylim(0, 1.1) # Beri ruang untuk label angka
    plt.grid(axis='y', linestyle='--', alpha=0.6)

    # Tandai Pemenang
    winner_idx = np.argmax(probs)
    plt.annotate('Kandidat Terpilih\n(Energi Terendah)', 
                 xy=(winner_idx, probs[winner_idx]), 
                 xytext=(winner_idx, probs[winner_idx] + 0.15),
                 arrowprops=dict(facecolor='black', shrink=0.05),
                 ha='center', fontweight='bold')

    filename = "visualisasi_probabilitas_vqe.png"
    plt.savefig(filename, dpi=300)
    plt.close()
    print(f"Grafik probabilitas berhasil disimpan sebagai: {filename}")

if __name__ == "__main__":
    # Gunakan parameter yang sesuai dengan run terakhir Anda
    render_vqe_probability_bar(n_qubits=2, K=1)
\end{lstlisting}

\begin{lstlisting}[language=Python, caption={Skrip visualisasi grafik interaksi antar aset (\texttt{project\_graph.py})}, label={lst:project-graph}]
import marimo

__generated_with = "0.23.5"
app = marimo.App()


@app.cell
def _():
    import marimo as mo
    import os
    import re

    return mo, os, re


@app.cell
def _(mo, os, re):
    # Fungsi untuk mencari file apa saja yang di-import oleh sebuah file
    def scan_dependencies(file_path, project_files):
        found_deps = set()
        try:
            with open(file_path, 'r', encoding='utf-8') as f:
                content = f.read()
                # Mencari pola 'import nama_file' atau 'from nama_file import ...'
                for other_file in project_files:
                    module_name = other_file.replace('.py', '')
                    # Regex untuk memastikan kita tidak salah mencocokkan kata di dalam string
                    pattern = rf'\b(import|from)\s+{module_name}\b'
                    if re.search(pattern, content):
                        found_deps.add(other_file)
        except:
            pass
        return found_deps

    def draw_project_map():
        root = os.getcwd()
        # Filter: Hanya ambil file .py, abaikan project_graph.py dan main_N...py
        all_py_files = [
            f for f in os.listdir(root) 
            if f.endswith('.py') 
            and f != 'project_graph.py'
            and not re.match(r'main_N\d+\.py', f)
        ]

        mermaid_lines = ["graph TD"]

        # Tambahkan gaya visual yang lebih kaya
        mermaid_lines.append("    classDef core fill:#FF9800,stroke:#E65100,stroke-width:2px,color:white;")
        mermaid_lines.append("    classDef strategy fill:#2196F3,stroke:#0D47A1,stroke-width:2px,color:white;")
        mermaid_lines.append("    classDef engine fill:#4CAF50,stroke:#1B5E20,stroke-width:2px,color:white;")
        mermaid_lines.append("    classDef math fill:#9C27B0,stroke:#4A148C,stroke-width:1px,color:white;")

        # Kategorisasi File
        categories = {
            'core': ['main.py', 'config.py', 'backtest_runner.py', 'report_generator.py', 'plot_generator.py', 'data_downloader.py'],
            'strategy': ['run_strategy_step.py', 'rebalance_portfolio.py'],
            'engine': ['run_vqe_adaptive.py', 'find_nash_sbr.py', 'brute_force_validator.py', 'find_optimal_lr_spsa.py', 'run_spsa_test.py'],
            'math': [f for f in all_py_files if f.startswith(('calc_', 'compute_', 'build_'))]
        }

        def get_style(filename):
            for cat, files in categories.items():
                if filename in files:
                    return f":::{cat}"
            return ""

        has_connections = False
        seen_edges = set()

        for file in all_py_files:
            deps = scan_dependencies(os.path.join(root, file), all_py_files)

            # Styling node berdasarkan kategori
            style = get_style(file)
            mermaid_lines.append(f'    {file.replace(".", "_")}["{file}"]{style}')

            for dep in deps:
                if file != dep:
                    src = file.replace(".", "_")
                    dst = dep.replace(".", "_")
                    edge = (src, dst)
                    if edge not in seen_edges:
                        mermaid_lines.append(f"    {src} --> {dst}")
                        seen_edges.add(edge)
                        has_connections = True
        if not has_connections:
            return mo.md("### ⚠️ Tidak ditemukan hubungan import antar file.\nPastikan file Anda menggunakan `import nama_file_lain`.")

        return mo.vstack([
            mo.md("# 🗺️ Peta Hubungan Kode Proyekmu"),
            mo.md("Garis panah menunjukkan file mana yang 'memanggil' file lainnya."),
            mo.mermaid("\n".join(mermaid_lines))
        ])

    draw_project_map()
    return


if __name__ == "__main__":
    app.run()
\end{lstlisting}

\begin{lstlisting}[language=Python, caption={Skrip utilitas untuk pengeksporan hasil graf fisis (\texttt{export\_graph.py})}, label={lst:export-graph}]
import os
import re
import base64
import requests # Pastikan library requests tersedia

def scan_dependencies(file_path, project_files):
    found_deps = set()
    try:
        with open(file_path, 'r', encoding='utf-8') as f:
            content = f.read()
            for other_file in project_files:
                module_name = other_file.replace('.py', '')
                pattern = rf'\b(import|from)\s+{module_name}\b'
                if re.search(pattern, content):
                    found_deps.add(other_file)
    except:
        pass
    return found_deps

def generate_mermaid():
    root = os.getcwd()
    all_py_files = [
        f for f in os.listdir(root) 
        if f.endswith('.py') 
        and f not in ['project_graph.py', 'export_graph.py']
        and not re.match(r'main_N\d+\.py', f)
    ]

    mermaid_lines = ["graph TD"]
    mermaid_lines.append("    classDef core fill:#FF9800,stroke:#E65100,stroke-width:2px,color:white;")
    mermaid_lines.append("    classDef strategy fill:#2196F3,stroke:#0D47A1,stroke-width:2px,color:white;")
    mermaid_lines.append("    classDef engine fill:#4CAF50,stroke:#1B5E20,stroke-width:2px,color:white;")
    mermaid_lines.append("    classDef math fill:#9C27B0,stroke:#4A148C,stroke-width:1px,color:white;")

    categories = {
        'core': ['main.py', 'config.py', 'backtest_runner.py', 'report_generator.py', 'plot_generator.py', 'data_downloader.py'],
        'strategy': ['run_strategy_step.py', 'rebalance_portfolio.py'],
        'engine': ['run_vqe_adaptive.py', 'find_nash_sbr.py', 'brute_force_validator.py', 'find_optimal_lr_spsa.py', 'run_spsa_test.py'],
        'math': [f for f in all_py_files if f.startswith(('calc_', 'compute_', 'build_'))]
    }

    def get_style(filename):
        for cat, files in categories.items():
            if filename in files:
                return f":::{cat}"
        return ""

    seen_edges = set()
    for file in all_py_files:
        style = get_style(file)
        mermaid_lines.append(f'    {file.replace(".", "_")}["{file}"]{style}')
        deps = scan_dependencies(os.path.join(root, file), all_py_files)
        for dep in deps:
            if file != dep:
                src = file.replace(".", "_")
                dst = dep.replace(".", "_")
                if (src, dst) not in seen_edges:
                    mermaid_lines.append(f"    {src} --> {dst}")
                    seen_edges.add((src, dst))
    
    return "\n".join(mermaid_lines)

def export_png(mermaid_text, output_file="project_graph.png"):
    print("Mengonversi Mermaid ke PNG via mermaid.ink...")
    # Encode ke base64
    graphbytes = mermaid_text.encode("ascii")
    base64_bytes = base64.b64encode(graphbytes)
    base64_string = base64_bytes.decode("ascii")
    
    url = "https://mermaid.ink/img/" + base64_string
    
    try:
        response = requests.get(url)
        if response.status_code == 200:
            with open(output_file, "wb") as f:
                f.write(response.content)
            print(f"Berhasil! Gambar disimpan sebagai: {output_file}")
        else:
            print(f"Gagal mengunduh gambar. Status code: {response.status_code}")
    except Exception as e:
        print(f"Terjadi kesalahan: {e}")

if __name__ == "__main__":
    mermaid_str = generate_mermaid()
    export_png(mermaid_str)
\end{lstlisting}

